
@@M0@@ d'après (1.3.ii)).

Remarquons enfin que si @@M1@@, @@M2@@.

\textbf{Théorème 1.6.} \uline{L'application de} @@M3@@ \uline{dans} @@M4@@ \uline{qui}, \uline{à un morphisme de foncteurs}, \uline{associe sa valeur sur l'élément} @@M5@@ \uline{de} @@M6@@, \uline{est injective}; \uline{son image est le sous-groupe de} @@M7@@ \uline{formé des éléments} @@M8@@ \uline{tels que} (\uline{avec les notations de} 1.4. \uline{pour les} @@M9@@)~:
@@M10@@

Montrons d'abord l'injectivité. D'après le lemme 1.2, il suffit de voir qu'un morphisme de foncteurs @@M11@@ nul sur @@M12@@ est nul sur @@M13@@. Le diagramme commutatif~:
@@M14@@
montre que pour tout @@M15@@, @@M16@@.

Or @@M17@@~; donc @@M18@@, et par suite @@M19@@, pour tout @@M20@@~; comme @@M21@@ est séparé et complet, @@M22@@.

Soit maintenant @@M23@@ un élément quelconque de @@M24@@, et soit 
